\documentclass[11pt,a4paper]{article}

% --- Packages ---
\usepackage[T1]{fontenc}
\usepackage[utf8]{inputenc}
\usepackage{lmodern}
\usepackage[english]{babel}
\usepackage[a4paper,margin=1in]{geometry}
\usepackage{microtype}
\usepackage{titlesec}
\usepackage{enumitem}
\usepackage{abstract}
\usepackage{authblk}
\usepackage{hyperref}
\usepackage{url}
\usepackage{xcolor}
\usepackage{fancyhdr}
\usepackage{booktabs}
\usepackage{tikz}
\usepackage{etoolbox}

% --- Colors ---
\definecolor{linkcol}{HTML}{004488}

\hypersetup{
    colorlinks=true,
    linkcolor=linkcol,
    urlcolor=linkcol,
    citecolor=linkcol,
    pdftitle={Genetic Algorithm for Multi-Stop Delivery Route Optimization},
    pdfauthor={Research Topic 1}
}

% --- Header / footer ---
\pagestyle{fancy}
\fancyhf{}
\fancyhead[L]{\small\itshape Topic 1 \textbullet{} Genetic Algorithm \textbullet{} Route Optimization}
\fancyhead[R]{\small\thepage}
\renewcommand{\headrulewidth}{0.4pt}

% --- Section style ---
\titleformat{\section}{\normalfont\large\bfseries\scshape}{\thesection.}{0.75em}{}
\titleformat{\subsection}{\normalfont\normalsize\bfseries\itshape}{\thesubsection}{0.6em}{}
\titlespacing*{\section}{0pt}{1.6em}{0.7em}
\titlespacing*{\subsection}{0pt}{1.2em}{0.5em}

% --- Abstract style ---
\renewcommand{\abstractnamefont}{\normalfont\bfseries\scshape}
\renewcommand{\abstracttextfont}{\normalfont\small\itshape}

% --- List spacing ---
\setlist[itemize]{leftmargin=1.3em,topsep=0.2em,itemsep=0.15em,parsep=0pt}
\setlist[enumerate]{leftmargin=1.5em,topsep=0.2em,itemsep=0.2em,parsep=0pt}

% --- Title block ---
\title{\vspace{-0.6in}\bfseries\Large Genetic Algorithm for Multi-Stop Delivery Route Optimization\\[2pt]
\large\mdseries\itshape A Practical Research Topic on Combinatorial Optimization}

\author[1]{Topic 1 \textbullet{} Algorithms \& Optimization}
\affil[1]{Practical Algorithm Research Series \textbullet{} 2026}

\date{}

\begin{document}
\maketitle
\thispagestyle{fancy}

\begin{abstract}
\noindent
A driver or delivery vehicle must visit many locations while minimizing total travel distance, travel time, fuel cost, or late deliveries. This problem is closely related to the Traveling Salesman Problem (TSP) and the Vehicle Routing Problem (VRP). This document proposes a practical research project that investigates whether a genetic algorithm can efficiently find near-optimal multi-stop delivery routes under real-world constraints such as distance, delivery deadlines, and vehicle capacity. It outlines the background, the core research question and sub-questions, a five-step research plan, the expected outcome, and source references.
\end{abstract}

\noindent\rule{\textwidth}{0.4pt}

% --- Background ---
\section{Basic Background}

Multi-stop delivery routing is a classic optimization problem. A driver or delivery vehicle must visit many locations while minimizing total travel distance, travel time, fuel cost, or late deliveries. This problem is closely related to the \textbf{Traveling Salesman Problem (TSP)} and the \textbf{Vehicle Routing Problem (VRP)}.

A simple shortest-path algorithm such as Dijkstra can find the shortest path between two locations, but it does not directly solve the harder question: \textbf{what is the best order to visit many delivery stops?} As the number of stops increases, the number of possible routes grows extremely quickly. For example, 10 stops already create millions of possible orderings.

A \textbf{genetic algorithm} is useful because it does not need to check every possible route. Instead, it starts with many possible routes, evaluates them using a fitness function, keeps better routes, combines parts of good routes, adds small mutations, and gradually improves the solution.

% --- Research Question ---
\section{Research Question}

\noindent\textit{\textbf{Can a genetic algorithm efficiently find near-optimal multi-stop delivery routes under real-world constraints such as distance, delivery deadlines, and vehicle capacity?}}

\subsection*{Possible sub-questions}
\begin{enumerate}
    \item How close is the genetic algorithm's route to the best-known or baseline route?
    \item How does performance change as the number of delivery stops increases?
    \item How do constraints such as time windows or vehicle capacity affect route quality?
    \item Can a hybrid method using Dijkstra plus genetic algorithm improve performance?
\end{enumerate}

% --- Research Plan ---
\section{Brief Research Plan}

\begin{enumerate}
    \item \textbf{Create or collect delivery-location data}
    \begin{itemize}
        \item Use a small city map, school campus map, or simulated delivery points.
        \item Each location has coordinates, demand size, and optional delivery deadline.
    \end{itemize}

    \item \textbf{Build baseline algorithms}
    \begin{itemize}
        \item Nearest-neighbor greedy algorithm.
        \item Random route generation.
        \item Optional: exact brute-force search for very small cases.
    \end{itemize}

    \item \textbf{Implement the genetic algorithm}
    \begin{itemize}
        \item Each chromosome represents one possible delivery order.
        \item Fitness function measures total distance, lateness penalty, and capacity penalty.
        \item Use selection, crossover, and mutation to evolve routes.
    \end{itemize}

    \item \textbf{Optional hybrid design}
    \begin{itemize}
        \item Use Dijkstra to compute shortest-path distances between all pairs of delivery locations.
        \item Use genetic algorithm to optimize the visiting order.
    \end{itemize}

    \item \textbf{Evaluate performance}
    \begin{itemize}
        \item Compare route distance, runtime, number of late deliveries, and scalability.
        \item Test with 10, 20, 50, and 100 stops.
    \end{itemize}
\end{enumerate}

% --- Expected Outcome ---
\section{Expected Outcome}

The project should show that genetic algorithms are practical for large multi-stop delivery problems where exact search is too slow. The student can also build a visualization showing how the route improves over generations.

% --- References ---
\section*{References}
\addcontentsline{toc}{section}{References}

\begin{enumerate}[label={[\arabic*]},leftmargin=2.2em]
    \item Sabah Sadiq, \textit{The Traveling Salesman Problem: Optimizing Delivery Routes Using Genetic Algorithms}, SAS Global Forum Proceedings, 2012.\\
    \url{https://support.sas.com/resources/papers/proceedings12/161-2012.pdf}

    \item Ahmed Hussain et al., \textit{Genetic Algorithm for Traveling Salesman Problem with Modified Cycle Crossover Operator}, Computational Intelligence and Neuroscience, 2017.\\
    \url{https://pmc.ncbi.nlm.nih.gov/articles/PMC5676484/}

    \item K. Bryant, \textit{Genetic Algorithms and the Travelling Salesman Problem}, Harvey Mudd College Senior Thesis, 2000.\\
    \url{https://scholarship.claremont.edu/cgi/viewcontent.cgi?article=1129&context=hmc_theses}
\end{enumerate}

\end{document}
